\documentclass[UTF8,12pt,a4paper]{ctexart}

\usepackage{amsmath}
\usepackage{cases}
\usepackage{cite}
\usepackage{graphicx}
\usepackage{enumerate}
\usepackage{algorithm}
\usepackage{caption} %\caption*需要
\usepackage[noend]{algpseudocode} %algorithmicx
\makeatletter
\renewcommand{\fnum@algorithm}{\fname@algorithm}
\makeatother
\renewcommand{\algorithmicrequire}{\textbf{Input:}}
\renewcommand{\algorithmicensure}{\textbf{Output:}}
\usepackage[margin=1in]{geometry}
\geometry{a4paper}
\usepackage{fancyhdr}
\pagestyle{fancy}
\fancyhf{}
\usepackage{xcolor}
\usepackage{listings}
\lstset{
	basicstyle=\ttfamily,                                % 传统上，代码打字机字体好些
	breaklines,
	tabsize=4,
	extendedchars=false,
	columns=fixed,
	%numbers=left,                                        % 在左侧显示行号
	numbers=none,
	frame=none,                                          % 不显示背景边框
	backgroundcolor=\color[RGB]{245,245,244},            % 设定背景颜色
	keywordstyle=\color[RGB]{40,40,255},                 % 设定关键字颜色
	numberstyle=\footnotesize\color{darkgray},           % 设定行号格式
	commentstyle=\it\color[RGB]{0,96,96},                % 设置代码注释的格式
	stringstyle=\rmfamily\slshape\color[RGB]{128,0,0},   % 设置字符串格式
	showstringspaces=false,                              % 不显示字符串中的空格
	xleftmargin=1em,
	%xrightmargin=1em,
	%aboveskip=-0.5em
	language=c++,                                        % 设置语言
}

% title依据是作业还是实验报告，以及是第几次作业/实验报告，
% 比如第1次分治算法作业，则命名为"算法原理 HW1 分治算法";
% 第1次分治算法实验报告，则命名为"算法原理 Lab1 分治算法"
\title{算法原理 HW1 作业}
\author{
	姓名：\underline{冯峻}~~~~~~
	学号：\underline{523031910148}~~~~~~}
\date{\today}
\pagenumbering{arabic}

\begin{document}

% 如下两行请按实际情况修改
\fancyhead[L]{冯峻}
\fancyhead[C]{HW1}
\fancyfoot[C]{\thepage}

\maketitle
%\tableofcontents
%\newpage

% \emph{以下内容仅针对作业，按照作业题目编号，将题目誊写过来，然后依次予以回答}
\begin{enumerate}[1.]
	\item \textbf{验证$f(n) = 2n^3+5n^2+7n+9$的复杂度是否是$\Theta(n^3)$？如是，请根据同阶函数集合的定义
	找出$c_1, c_2, n_0$；如不是，请给出证据。}

	答：\textbf{其复杂度是$\Theta(n^3)$}，参数取值以及理由如下：
	
	设$c_1 n^3 \leq 2n^3+5n^2+7n+9 \leq c_2 n^3$

	易知，如果要使得左式成立，只需要$c_1 = 2$，则在$n > 0$时即可满足$c_1 n^3 \leq 2n^3+5n^2+7n+9$时恒成立

	现在我们来论证右式的成立性，我们不妨直接令$c_2 = 23$，那么右式等价为：
	\begin{equation}
		2n^3+5n^2+7n+9 \leq 2n^3+5n^3+7n^3+9n^3
		\label{eq:T1}
	\end{equation}

	易知该式在$n\geq1$时恒成立。

	因此，我们可以取：$c_1 = 2, c_2 = 23, n_0 = 1$
	\item \textbf{参考课件，求解递归方程$T(n)=2T(\sqrt{n})+3\text{log}n$}

	答：

	令$m=\text{log}_2 n$，带入原递归方程后递归方程化为：
	\begin{equation}
		T(2^m)=2T(2^{m/2})+3m
		\label{eq:T2_trans}
	\end{equation}
	记$S(m) = T(2^m)$，则方程\eqref{eq:T2_trans}等效为：
	\begin{equation}
		S(m) = S(m/2) + 3m
		\label{eq:T_trans_2}
	\end{equation}
	有$\Theta(m^{\text{log}_{b}a})=\Theta(m)\text{与}3m\text{同阶}$
	
	因此根据主定理，我们有：
	\begin{equation}
		S(m) = \Theta(m\text{log}m)
	\end{equation}
	所以有：
	\begin{equation}
		T(n) = S(\text{log}m) = \Theta(\log(m)\log(\log(m)))
	\end{equation}

	\item \textbf{求解递归方程$T(n)=4T(n/2)+n^2\text{log}^2n$的上界。注意：如果要用主定理，请严格按照课
	件判定适用性。}

	答：

	因为我们只需要求解该复杂度的上界，所以我们不妨构造这样一个递归方程：
	\begin{equation}
		S(n)=4S(n/2)+n^3
		\label{eq:T3_up}
	\end{equation}
	已知$\sqrt{n} > \text{log}n$对任意$n>0$恒成立，因此有：
	\begin{equation}
		n^3 > n^2\text{log}^2 n
		\label{eq:T3_n_log_relation}
	\end{equation}
	那么由递归方程易知：$T(n) < S(n)$，即$S(n)$是$T(n)$的一个上界

	因此，我们只需要求解递归方程\eqref{eq:T3_up}，就能求解得到$T(n)$的一个上界。

	根据主定理的形式：在递归方程\eqref{eq:T3_up}中，有$a=4, b=2, f(n) = n^3$

	根据主定理的第三条，对于大于0的$\epsilon$，例如取$\epsilon=0.5$，我们有：$f(n) = n^3 = \Omega(n^{\text{log}_b a + \epsilon}) = \Omega({n^{2+\epsilon}})$

	同时，对于$n>1$的任意$n$，取$c=0.6<1$，我们有：$af(n/b)=4(\frac{n}{2})^3=\frac{n^3}{2}<0.6n^3=cf(n)$，满足正则条件，因此有：
	\begin{equation}
		S(n) = \Theta(n^3)
		\label{eq:S}
	\end{equation}
	因此，我们求得了$T(n)$的一个上界：$\Theta({n^3})$

\end{enumerate}



\end{document}
